\section{Comparison and discussion}
\label{sec:discussion}

\subsection{Comparison to existing automation}
\label{sec:discussion-comparison}

\paragraph{aesop.}
Aesop \citep{aesop2023} is the workhorse forward-search tactic in
Lean 4. \texttt{pythia} declares a named aesop ruleset and uses
aesop in stage 2 of the dispatcher chain. \texttt{pythia} is not a
replacement for aesop. It is a statistics-domain registry plus a
Z3 oracle plus a vacuous-closure audit, layered on top of aesop.

\paragraph{Sledgehammer.}
Sledgehammer \citep{paulsonblanchette2010threeyears,
blanchettesledgehammer2011} dispatches Isabelle/HOL goals to external
ATPs (Vampire, E, Z3, CVC4) and reconstructs the proof. There is no
direct Lean 4 equivalent. \texttt{pythia} addresses the
statistics-domain case via \texttt{z3\_check} for linear-real-arithmetic and
Aristotle as a structural oracle. A general-purpose Lean
Sledgehammer is out of scope for this paper.

\paragraph{CoqHammer.}
CoqHammer \citep{czajkakaliszyk2018coqhammer} dispatches Coq goals
to external ATPs with kernel-checked reconstruction. \texttt{z3\_check}
in \texttt{pythia} adopts the same reconstruction template
(witness from external prover, in-kernel checking via
\texttt{linarith}). CoqHammer is the design template. We do not
claim a Lean-side reimplementation of CoqHammer's full premise
selection.

\paragraph{Mathlib's existing tactic suite.}
Mathlib supplies \texttt{linarith}, \texttt{polyrith},
\texttt{nlinarith}, \texttt{positivity}, \texttt{measurability},
\texttt{bound}, \texttt{gcongr}. None of these is a statistics-domain
hammer. \texttt{pythia} composes with all of them. The fall-through
to \texttt{linarith} in stage 3 of the dispatcher chain is the
explicit composition.

\subsection{Aristotle vs DSPv2 vs Opus 4.6 (illustrative)}
\label{sec:discussion-comparison-empirical}

% TODO: pull 3-4 problems from FormalAVS where Aristotle closes
% and hand-Mathlib alone does not. asabi to populate from
% neurips-2026-dataset-paper tables.

Table~\ref{tab:illustrative-comparison} reports a small subset of
FormalAVS problems \citep{anonymous2026formalavs} on which
Aristotle, DSPv2 (a fine-tuned local-GPU drafter), and Opus 4.6
(an API model) were each given a single attempt. The full benchmark
numbers are in the dataset paper. The point of this subset is to
illustrate the regime in which Aristotle as oracle closes problems
that hand-Mathlib alone does not. Per-problem solve outcomes (close
or open) are kernel-checked.

\begin{table}[h]
\centering
\small
\begin{tabular}{lccc}
\toprule
\textbf{Problem} & \textbf{Aristotle} & \textbf{DSPv2} & \textbf{Opus 4.6} \\
\midrule
% TODO: replace placeholders with real solve/no-solve outcomes
% from the FormalAVS subset chosen for this illustrative comparison.
combine\_eprocesses\_avg & TODO & TODO & TODO \\
ville\_bound\_betting & TODO & TODO & TODO \\
matrix\_bernstein\_step & TODO & TODO & TODO \\
bretagnolle\_huber\_binary & TODO & TODO & TODO \\
\bottomrule
\end{tabular}
\caption{Illustrative subset of FormalAVS. Solve outcomes
(close vs open) per single attempt. Full benchmark numbers in
\citet{anonymous2026formalavs}.}
\label{tab:illustrative-comparison}
\end{table}

The wider table on the full benchmark subset is in
Appendix~\ref{sec:appendix-comparison}.

\subsection{Limitations}
\label{sec:discussion-limitations}

Three concrete gaps. (i) Domain scope: \texttt{pythia} is
statistics-domain only. It does not address general proof automation
in Lean 4. (ii) Toolchain pin: the library ships pinned to Lean
4.28 + Mathlib v4.28.0. Mathlib evolves quickly. Each major Mathlib
release typically requires a porting pass. (iii) Aristotle
dependency: the hardest closures route to a proprietary external
prover. The \texttt{z3\_check} oracle is open-source. The Aristotle
oracle is not. The library is usable without Aristotle for any goal
that closes via the \texttt{pythia} dispatcher chain or
\texttt{z3\_check}. A subset of structural Mathlib gaps requires
either Aristotle or a hand proof.
